%% Appendix "An Introduction to the Framework": Section 2 of the vqupit
%% paper (vqupit/sections/permutations.lagda.tex), the swap circuits
%% presenting the symmetric groups.  make.sh writes
%% ../latex/pldi/lagda/app-permutations.tex; edit this file.
\begin{code}[hide]%
\>[0]\AgdaKeyword{module}\AgdaSpace{}%
\AgdaModule{pldi.lagda.app-permutations}\AgdaSpace{}%
\AgdaKeyword{where}\<%
\\
%
\\[\AgdaEmptyExtraSkip]%
\>[0]\AgdaKeyword{open}\AgdaSpace{}%
\AgdaKeyword{import}\AgdaSpace{}%
\AgdaModule{Data.Nat}\AgdaSpace{}%
\AgdaKeyword{using}\AgdaSpace{}%
\AgdaSymbol{(}\AgdaDatatype{ℕ}\AgdaSpace{}%
\AgdaSymbol{;}\AgdaSpace{}%
\AgdaInductiveConstructor{suc}\AgdaSymbol{)}\<%
\\
\>[0]\AgdaKeyword{open}\AgdaSpace{}%
\AgdaKeyword{import}\AgdaSpace{}%
\AgdaModule{Data.Product}\AgdaSpace{}%
\AgdaKeyword{using}\AgdaSpace{}%
\AgdaSymbol{(}\AgdaOperator{\AgdaFunction{\AgdaUnderscore{}×\AgdaUnderscore{}}}\AgdaSpace{}%
\AgdaSymbol{;}\AgdaSpace{}%
\AgdaOperator{\AgdaInductiveConstructor{\AgdaUnderscore{},\AgdaUnderscore{}}}\AgdaSpace{}%
\AgdaSymbol{;}\AgdaSpace{}%
\AgdaField{proj₁}\AgdaSpace{}%
\AgdaSymbol{;}\AgdaSpace{}%
\AgdaField{proj₂}\AgdaSymbol{)}\<%
\\
\>[0]\AgdaKeyword{open}\AgdaSpace{}%
\AgdaKeyword{import}\AgdaSpace{}%
\AgdaModule{Data.Unit}\AgdaSpace{}%
\AgdaKeyword{using}\AgdaSpace{}%
\AgdaSymbol{(}\AgdaRecord{⊤}\AgdaSymbol{)}\<%
\\
\>[0]\AgdaKeyword{open}\AgdaSpace{}%
\AgdaKeyword{import}\AgdaSpace{}%
\AgdaModule{Level}\AgdaSpace{}%
\AgdaKeyword{using}\AgdaSpace{}%
\AgdaSymbol{(}\AgdaFunction{0ℓ}\AgdaSpace{}%
\AgdaSymbol{;}\AgdaSpace{}%
\AgdaPostulate{Level}\AgdaSpace{}%
\AgdaSymbol{;}\AgdaSpace{}%
\AgdaOperator{\AgdaPrimitive{\AgdaUnderscore{}⊔\AgdaUnderscore{}}}\AgdaSymbol{)}\<%
\\
\>[0]\AgdaKeyword{open}\AgdaSpace{}%
\AgdaKeyword{import}\AgdaSpace{}%
\AgdaModule{Relation.Binary}\AgdaSpace{}%
\AgdaKeyword{using}\AgdaSpace{}%
\AgdaSymbol{(}\AgdaFunction{Rel}\AgdaSymbol{)}\<%
\\
\>[0]\AgdaKeyword{open}\AgdaSpace{}%
\AgdaKeyword{import}\AgdaSpace{}%
\AgdaModule{Algebra.Bundles}\AgdaSpace{}%
\AgdaKeyword{using}\AgdaSpace{}%
\AgdaSymbol{(}\AgdaRecord{Group}\AgdaSymbol{)}\<%
\\
%
\\[\AgdaEmptyExtraSkip]%
\>[0]\AgdaKeyword{open}\AgdaSpace{}%
\AgdaKeyword{import}\AgdaSpace{}%
\AgdaModule{Notations}\AgdaSpace{}%
\AgdaKeyword{using}\AgdaSpace{}%
\AgdaSymbol{(}\AgdaInductiveConstructor{₁₊}\AgdaSpace{}%
\AgdaSymbol{;}\AgdaSpace{}%
\AgdaInductiveConstructor{₂₊}\AgdaSpace{}%
\AgdaSymbol{;}\AgdaSpace{}%
\AgdaInductiveConstructor{₃₊}\AgdaSymbol{)}\<%
\\
\>[0]\AgdaKeyword{open}\AgdaSpace{}%
\AgdaKeyword{import}\AgdaSpace{}%
\AgdaModule{Word.Base}\AgdaSpace{}%
\AgdaKeyword{using}\AgdaSpace{}%
\AgdaSymbol{(}\AgdaDatatype{Word}\AgdaSpace{}%
\AgdaSymbol{;}\AgdaSpace{}%
\AgdaFunction{WRel}\AgdaSpace{}%
\AgdaSymbol{;}\AgdaSpace{}%
\AgdaOperator{\AgdaInductiveConstructor{[\AgdaUnderscore{}]ʷ}}\AgdaSpace{}%
\AgdaSymbol{;}\AgdaSpace{}%
\AgdaInductiveConstructor{ε}\AgdaSpace{}%
\AgdaSymbol{;}\AgdaSpace{}%
\AgdaOperator{\AgdaInductiveConstructor{\AgdaUnderscore{}•\AgdaUnderscore{}}}\AgdaSpace{}%
\AgdaSymbol{;}\AgdaSpace{}%
\AgdaFunction{wmap}\AgdaSymbol{)}\<%
\\
\>[0]\AgdaKeyword{open}\AgdaSpace{}%
\AgdaKeyword{import}\AgdaSpace{}%
\AgdaModule{Presentation.GroupLike}\AgdaSpace{}%
\AgdaKeyword{using}\AgdaSpace{}%
\AgdaSymbol{(}\AgdaFunction{Grouplike}\AgdaSymbol{)}\<%
\\
\>[0]\AgdaKeyword{open}\AgdaSpace{}%
\AgdaKeyword{import}\AgdaSpace{}%
\AgdaModule{Examples.Groups.Symmetric.Semantics}\AgdaSpace{}%
\AgdaKeyword{using}\AgdaSpace{}%
\AgdaSymbol{(}\AgdaFunction{Permutation′-group}\AgdaSymbol{)}\<%
\\
\>[0]\AgdaKeyword{import}\AgdaSpace{}%
\AgdaModule{Examples.Groups.Symmetric.Presentation}\AgdaSpace{}%
\AgdaSymbol{as}\AgdaSpace{}%
\AgdaModule{SymPres}\<%
\\
%
\\[\AgdaEmptyExtraSkip]%
\>[0]\AgdaKeyword{private}\AgdaSpace{}%
\AgdaKeyword{variable}\<%
\\
\>[0][@{}l@{\AgdaIndent{0}}]%
\>[2]\AgdaGeneralizable{n}\AgdaSpace{}%
\AgdaSymbol{:}\AgdaSpace{}%
\AgdaDatatype{ℕ}\<%
\\
%
\>[2]\AgdaGeneralizable{a}\AgdaSpace{}%
\AgdaGeneralizable{ℓ}\AgdaSpace{}%
\AgdaSymbol{:}\AgdaSpace{}%
\AgdaPostulate{Level}\<%
\\
%
\>[2]\AgdaGeneralizable{X}\AgdaSpace{}%
\AgdaSymbol{:}\AgdaSpace{}%
\AgdaPrimitive{Set}\<%
\\
%
\\[\AgdaEmptyExtraSkip]%
\>[0]\AgdaKeyword{postulate}\<%
\\
\>[0][@{}l@{\AgdaIndent{0}}]%
\>[2]\AgdaPostulate{proof-elided}\AgdaSpace{}%
\AgdaSymbol{:}\AgdaSpace{}%
\AgdaSymbol{∀}\AgdaSpace{}%
\AgdaSymbol{\{}\AgdaBound{ℓ}\AgdaSymbol{\}}\AgdaSpace{}%
\AgdaSymbol{\{}\AgdaBound{A}\AgdaSpace{}%
\AgdaSymbol{:}\AgdaSpace{}%
\AgdaPrimitive{Set}\AgdaSpace{}%
\AgdaBound{ℓ}\AgdaSymbol{\}}\AgdaSpace{}%
\AgdaSymbol{→}\AgdaSpace{}%
\AgdaBound{A}\<%
\\
%
\>[2]\AgdaPostulate{IsGroupIsomorphism}\AgdaSpace{}%
\AgdaSymbol{:}\AgdaSpace{}%
\AgdaSymbol{\{}\AgdaBound{G}\AgdaSpace{}%
\AgdaSymbol{:}\AgdaSpace{}%
\AgdaRecord{Group}\AgdaSpace{}%
\AgdaGeneralizable{a}\AgdaSpace{}%
\AgdaGeneralizable{ℓ}\AgdaSymbol{\}}\AgdaSpace{}%
\AgdaSymbol{→}\AgdaSpace{}%
\AgdaSymbol{(}\AgdaDatatype{Word}\AgdaSpace{}%
\AgdaGeneralizable{X}\AgdaSpace{}%
\AgdaSymbol{→}\AgdaSpace{}%
\AgdaField{Group.Carrier}\AgdaSpace{}%
\AgdaBound{G}\AgdaSymbol{)}\AgdaSpace{}%
\AgdaSymbol{→}\AgdaSpace{}%
\AgdaPrimitive{Set}\AgdaSpace{}%
\AgdaSymbol{(}\AgdaGeneralizable{a}\AgdaSpace{}%
\AgdaOperator{\AgdaPrimitive{⊔}}\AgdaSpace{}%
\AgdaGeneralizable{ℓ}\AgdaSymbol{)}\<%
\end{code}
\section{The EquaCert Framework}\label{sec:form}

EquaCert is an extension of an existing library for verifying
completeness of circuit equational theory~\cite{bian2027popl}. We
first introduce the basics through an exmaple and then explain the
extension. The example is presenting symmetric group through circuits
of swap gate, which shows how gates, relations, cosets, tower,
semantics, uniqueness, completeness, presentation are defined in code.
\subsection{Gates, Generators, and Circuits}\label{app:sn-gates}

A gate set is a family $\aid{Gate} : \mathbb{N} \to \mathsf{Set}$
giving, for each arity, the gates of that arity; for permutation
circuit there is one gate, on two wires. A generic wraper $\aid{Gen}$
turns any such family into wire-indexed \emph{generators}: a gate
placed at the bottom of the wires, or a generator shifted up by one
wire.

\begin{code}%
\>[0]\AgdaKeyword{data}\AgdaSpace{}%
\AgdaDatatype{Gate}\AgdaSpace{}%
\AgdaSymbol{:}\AgdaSpace{}%
\AgdaDatatype{ℕ}\AgdaSpace{}%
\AgdaSymbol{→}\AgdaSpace{}%
\AgdaPrimitive{Set}\AgdaSpace{}%
\AgdaKeyword{where}\<%
\\
\>[0][@{}l@{\AgdaIndent{0}}]%
\>[2]\AgdaInductiveConstructor{σ-gate}\AgdaSpace{}%
\AgdaSymbol{:}\AgdaSpace{}%
\AgdaDatatype{Gate}\AgdaSpace{}%
\AgdaNumber{2}\<%
\\
\>[0]\AgdaKeyword{data}\AgdaSpace{}%
\AgdaDatatype{Gen}\AgdaSpace{}%
\AgdaSymbol{:}\AgdaSpace{}%
\AgdaDatatype{ℕ}\AgdaSpace{}%
\AgdaSymbol{→}\AgdaSpace{}%
\AgdaPrimitive{Set}\AgdaSpace{}%
\AgdaKeyword{where}\<%
\\
\>[0][@{}l@{\AgdaIndent{0}}]%
\>[2]\AgdaInductiveConstructor{gate₁}\AgdaSpace{}%
\AgdaSymbol{:}\AgdaSpace{}%
\AgdaDatatype{Gate}\AgdaSpace{}%
\AgdaNumber{1}\AgdaSpace{}%
\AgdaSymbol{→}\AgdaSpace{}%
\AgdaDatatype{Gen}\AgdaSpace{}%
\AgdaSymbol{(}\AgdaInductiveConstructor{₁₊}\AgdaSpace{}%
\AgdaGeneralizable{n}\AgdaSymbol{)}%
\>[34]\AgdaComment{--\ a\ 1-ary\ gate\ on\ the\ bottom\ wire}\<%
\\
%
\>[2]\AgdaInductiveConstructor{gate₂}\AgdaSpace{}%
\AgdaSymbol{:}\AgdaSpace{}%
\AgdaDatatype{Gate}\AgdaSpace{}%
\AgdaNumber{2}\AgdaSpace{}%
\AgdaSymbol{→}\AgdaSpace{}%
\AgdaDatatype{Gen}\AgdaSpace{}%
\AgdaSymbol{(}\AgdaInductiveConstructor{₂₊}\AgdaSpace{}%
\AgdaGeneralizable{n}\AgdaSymbol{)}%
\>[34]\AgdaComment{--\ a\ 2-ary\ gate\ on\ the\ bottom\ two\ wires}\<%
\\
%
\>[2]\AgdaOperator{\AgdaInductiveConstructor{\AgdaUnderscore{}↥}}%
\>[7]\AgdaSymbol{:}\AgdaSpace{}%
\AgdaDatatype{Gen}\AgdaSpace{}%
\AgdaGeneralizable{n}\AgdaSpace{}%
\AgdaSymbol{→}\AgdaSpace{}%
\AgdaDatatype{Gen}\AgdaSpace{}%
\AgdaSymbol{(}\AgdaInductiveConstructor{₁₊}\AgdaSpace{}%
\AgdaGeneralizable{n}\AgdaSymbol{)}%
\>[34]\AgdaComment{--\ the\ same\ generator,\ one\ wire\ higher}\<%
\end{code}
\begin{code}[hide]%
\>[0]\AgdaFunction{Circuit}\AgdaSpace{}%
\AgdaSymbol{:}\AgdaSpace{}%
\AgdaDatatype{ℕ}\AgdaSpace{}%
\AgdaSymbol{→}\AgdaSpace{}%
\AgdaPrimitive{Set}\<%
\\
\>[0]\AgdaFunction{Circuit}\AgdaSpace{}%
\AgdaBound{n}\AgdaSpace{}%
\AgdaSymbol{=}\AgdaSpace{}%
\AgdaDatatype{Word}\AgdaSpace{}%
\AgdaSymbol{(}\AgdaDatatype{Gen}\AgdaSpace{}%
\AgdaBound{n}\AgdaSymbol{)}\<%
\\
%
\\[\AgdaEmptyExtraSkip]%
\>[0]\AgdaFunction{CRel}\AgdaSpace{}%
\AgdaSymbol{:}\AgdaSpace{}%
\AgdaDatatype{ℕ}\AgdaSpace{}%
\AgdaSymbol{→}\AgdaSpace{}%
\AgdaPrimitive{Set₁}\<%
\\
\>[0]\AgdaFunction{CRel}\AgdaSpace{}%
\AgdaBound{n}\AgdaSpace{}%
\AgdaSymbol{=}\AgdaSpace{}%
\AgdaFunction{Rel}\AgdaSpace{}%
\AgdaSymbol{(}\AgdaFunction{Circuit}\AgdaSpace{}%
\AgdaBound{n}\AgdaSymbol{)}\AgdaSpace{}%
\AgdaFunction{0ℓ}\<%
\\
%
\\[\AgdaEmptyExtraSkip]%
\>[0]\AgdaOperator{\AgdaFunction{\AgdaUnderscore{}↑}}\AgdaSpace{}%
\AgdaSymbol{:}\AgdaSpace{}%
\AgdaFunction{Circuit}\AgdaSpace{}%
\AgdaGeneralizable{n}\AgdaSpace{}%
\AgdaSymbol{→}\AgdaSpace{}%
\AgdaFunction{Circuit}\AgdaSpace{}%
\AgdaSymbol{(}\AgdaInductiveConstructor{₁₊}\AgdaSpace{}%
\AgdaGeneralizable{n}\AgdaSymbol{)}\<%
\\
\>[0]\AgdaOperator{\AgdaFunction{\AgdaUnderscore{}↑}}\AgdaSpace{}%
\AgdaSymbol{=}\AgdaSpace{}%
\AgdaFunction{wmap}\AgdaSpace{}%
\AgdaOperator{\AgdaInductiveConstructor{\AgdaUnderscore{}↥}}\<%
\\
%
\\[\AgdaEmptyExtraSkip]%
\>[0]\AgdaKeyword{pattern}\AgdaSpace{}%
\AgdaInductiveConstructor{σ-gen}\AgdaSpace{}%
\AgdaSymbol{=}\AgdaSpace{}%
\AgdaInductiveConstructor{gate₂}\AgdaSpace{}%
\AgdaInductiveConstructor{σ-gate}\<%
\\
%
\\[\AgdaEmptyExtraSkip]%
\>[0]\AgdaFunction{σ}\AgdaSpace{}%
\AgdaSymbol{:}\AgdaSpace{}%
\AgdaFunction{Circuit}\AgdaSpace{}%
\AgdaSymbol{(}\AgdaInductiveConstructor{₂₊}\AgdaSpace{}%
\AgdaGeneralizable{n}\AgdaSymbol{)}\<%
\\
\>[0]\AgdaFunction{σ}\AgdaSpace{}%
\AgdaSymbol{=}\AgdaSpace{}%
\AgdaOperator{\AgdaInductiveConstructor{[}}\AgdaSpace{}%
\AgdaInductiveConstructor{σ-gen}\AgdaSpace{}%
\AgdaOperator{\AgdaInductiveConstructor{]ʷ}}\<%
\\
%
\\[\AgdaEmptyExtraSkip]%
\>[0]\AgdaKeyword{infix}\AgdaSpace{}%
\AgdaNumber{4}\AgdaSpace{}%
\AgdaOperator{\AgdaDatatype{\AgdaUnderscore{}SRel,\AgdaUnderscore{}===\AgdaUnderscore{}}}\AgdaSpace{}%
\AgdaOperator{\AgdaDatatype{\AgdaUnderscore{}VRel,\AgdaUnderscore{}===\AgdaUnderscore{}}}\AgdaSpace{}%
\AgdaOperator{\AgdaPostulate{\AgdaUnderscore{}≈\AgdaUnderscore{}}}\<%
\\
\>[0]\AgdaKeyword{infixr}\AgdaSpace{}%
\AgdaNumber{10}\AgdaSpace{}%
\AgdaOperator{\AgdaInductiveConstructor{σ•\AgdaUnderscore{}}}\<%
\\
%
\\[\AgdaEmptyExtraSkip]%
\>[0]\AgdaKeyword{private}\AgdaSpace{}%
\AgdaKeyword{variable}\<%
\\
\>[0][@{}l@{\AgdaIndent{0}}]%
\>[2]\AgdaGeneralizable{w}\AgdaSpace{}%
\AgdaGeneralizable{v}\AgdaSpace{}%
\AgdaSymbol{:}\AgdaSpace{}%
\AgdaFunction{Circuit}\AgdaSpace{}%
\AgdaGeneralizable{n}\<%
\\
%
\\[\AgdaEmptyExtraSkip]%
\>[0]\AgdaKeyword{postulate}\<%
\\
\>[0][@{}l@{\AgdaIndent{0}}]%
\>[2]\AgdaOperator{\AgdaPostulate{\AgdaUnderscore{}≈\AgdaUnderscore{}}}\AgdaSpace{}%
\AgdaSymbol{:}\AgdaSpace{}%
\AgdaFunction{Circuit}\AgdaSpace{}%
\AgdaGeneralizable{n}\AgdaSpace{}%
\AgdaSymbol{→}\AgdaSpace{}%
\AgdaFunction{Circuit}\AgdaSpace{}%
\AgdaGeneralizable{n}\AgdaSpace{}%
\AgdaSymbol{→}\AgdaSpace{}%
\AgdaPrimitive{Set}\<%
\end{code}

Here $n$ is a uninstantiated parameter, so $\aid{gate₂}\;\aid{σ-gate}$
is a generator at any width from two upward.  An $n$-wire
\emph{circuit} is a word of generators,
$\aid{Circuit}\;n = \aid{Word}\;(\aid{Gen}\;n)$, which is a single
generator $[\,g\,]$\aid{ʷ}, or the empty circuit \aid{ε} ($n$ bare
wires, the identity circuit at that width), or a sequential
composition $c \aidop{•} c'$.  We write $\sigma$ for the one-letter
circuit $[\,\aid{gate₂}\;\aid{σ-gate}\;]$\aid{ʷ} and $c\;{\uparrow}$
for $c$ shifted up one wire. Read diagrammatically, with wires drawn
bottom-up and a word drawn left to right as written, $\sigma$ is a
crossing of the bottom two wires, $\sigma\;{\uparrow}$ the same
crossing one wire higher, and
$\sigma \aidop{•} \sigma\;{\uparrow} \aidop{•} \sigma$ the three-wire
braid on the left of the \aid{yang-baxter} equation drawn beside its
axiom below.

\subsection{Relations and structural
  relations}\label{app:sn-relations}

The relations are a familiy of pairs indexed by the wire count --- the
uninstantiated $n$. Simlilarly to circuit, one relation covers every
width at which the relation fits. Agda names may be mixfix: each
underscore in \aid{{\_}SRel,{\_}==={\_}} marks an argument position,
so the family is applied as $n\;\aid{SRel,}\;w \aidrel{===} v$, read
``$w = v$ is an axiom on $n$ wires''.

\begin{center}
\begin{minipage}[c]{0.58\linewidth}
\begingroup\setlength{\mathindent}{0pt}
\begin{code}%
\>[0]\AgdaKeyword{data}\AgdaSpace{}%
\AgdaOperator{\AgdaDatatype{\AgdaUnderscore{}SRel,\AgdaUnderscore{}===\AgdaUnderscore{}}}\AgdaSpace{}%
\AgdaSymbol{:}\AgdaSpace{}%
\AgdaSymbol{(}\AgdaBound{n}\AgdaSpace{}%
\AgdaSymbol{:}\AgdaSpace{}%
\AgdaDatatype{ℕ}\AgdaSymbol{)}\AgdaSpace{}%
\AgdaSymbol{→}\AgdaSpace{}%
\AgdaFunction{WRel}\AgdaSpace{}%
\AgdaSymbol{(}\AgdaDatatype{Gen}\AgdaSpace{}%
\AgdaBound{n}\AgdaSymbol{)}\AgdaSpace{}%
\AgdaKeyword{where}\<%
\\
\>[0][@{}l@{\AgdaIndent{0}}]%
\>[2]\AgdaInductiveConstructor{order}%
\>[14]\AgdaSymbol{:}\AgdaSpace{}%
\AgdaSymbol{(}\AgdaInductiveConstructor{₂₊}\AgdaSpace{}%
\AgdaGeneralizable{n}\AgdaSymbol{)}\AgdaSpace{}%
\AgdaOperator{\AgdaDatatype{SRel,}}\AgdaSpace{}%
\AgdaFunction{σ}\AgdaSpace{}%
\AgdaOperator{\AgdaInductiveConstructor{•}}\AgdaSpace{}%
\AgdaFunction{σ}\AgdaSpace{}%
\AgdaOperator{\AgdaDatatype{===}}\AgdaSpace{}%
\AgdaInductiveConstructor{ε}\<%
\\
%
\>[2]\AgdaInductiveConstructor{yang-baxter}\AgdaSpace{}%
\AgdaSymbol{:}\AgdaSpace{}%
\AgdaSymbol{(}\AgdaInductiveConstructor{₃₊}\AgdaSpace{}%
\AgdaGeneralizable{n}\AgdaSymbol{)}\AgdaSpace{}%
\AgdaOperator{\AgdaDatatype{SRel,}}\AgdaSpace{}%
\AgdaFunction{σ}\AgdaSpace{}%
\AgdaOperator{\AgdaInductiveConstructor{•}}\AgdaSpace{}%
\AgdaFunction{σ}\AgdaSpace{}%
\AgdaOperator{\AgdaFunction{↑}}\AgdaSpace{}%
\AgdaOperator{\AgdaInductiveConstructor{•}}\AgdaSpace{}%
\AgdaFunction{σ}\AgdaSpace{}%
\AgdaOperator{\AgdaDatatype{===}}\AgdaSpace{}%
\AgdaFunction{σ}\AgdaSpace{}%
\AgdaOperator{\AgdaFunction{↑}}\AgdaSpace{}%
\AgdaOperator{\AgdaInductiveConstructor{•}}\AgdaSpace{}%
\AgdaFunction{σ}\AgdaSpace{}%
\AgdaOperator{\AgdaInductiveConstructor{•}}\AgdaSpace{}%
\AgdaFunction{σ}\AgdaSpace{}%
\AgdaOperator{\AgdaFunction{↑}}\<%
\end{code}
\endgroup
\end{minipage}\hspace{2.5em}%
\begin{minipage}[c]{0.28\linewidth}\centering
  \snfig[0.35]{sn-order}\\[4pt]
  \snfig[0.35]{sn-yb}
\end{minipage}
\end{center}

\noindent These two axioms are two thirds of the Coxeter presentation
of $S_n$: $\sigma^2 = \varepsilon$ and the braid relation.  The third
Coxeter relation that swaps on disjoint wires commute, is not specific
to permutations. It is one of the relations every circuit theory
needs, whatever its gates.  In code, \aid{Lift-Relation} extends any
relation family with the structural relations to the full relation set
\aid{{\_}VRel,{\_}==={\_}}:

\begin{code}%
\>[0]\AgdaKeyword{data}\AgdaSpace{}%
\AgdaOperator{\AgdaDatatype{\AgdaUnderscore{}VRel,\AgdaUnderscore{}===\AgdaUnderscore{}}}\AgdaSpace{}%
\AgdaSymbol{:}\AgdaSpace{}%
\AgdaSymbol{(}\AgdaBound{n}\AgdaSpace{}%
\AgdaSymbol{:}\AgdaSpace{}%
\AgdaDatatype{ℕ}\AgdaSymbol{)}\AgdaSpace{}%
\AgdaSymbol{→}\AgdaSpace{}%
\AgdaFunction{CRel}\AgdaSpace{}%
\AgdaBound{n}\AgdaSpace{}%
\AgdaKeyword{where}\<%
\\
\>[0][@{}l@{\AgdaIndent{0}}]%
\>[2]\AgdaInductiveConstructor{srel}%
\>[8]\AgdaSymbol{:}\AgdaSpace{}%
\AgdaGeneralizable{n}\AgdaSpace{}%
\AgdaOperator{\AgdaDatatype{SRel,}}\AgdaSpace{}%
\AgdaGeneralizable{w}\AgdaSpace{}%
\AgdaOperator{\AgdaDatatype{===}}\AgdaSpace{}%
\AgdaGeneralizable{v}\AgdaSpace{}%
\AgdaSymbol{→}\AgdaSpace{}%
\AgdaGeneralizable{n}\AgdaSpace{}%
\AgdaOperator{\AgdaDatatype{VRel,}}\AgdaSpace{}%
\AgdaGeneralizable{w}\AgdaSpace{}%
\AgdaOperator{\AgdaDatatype{===}}\AgdaSpace{}%
\AgdaGeneralizable{v}\<%
\\
%
\>[2]\AgdaInductiveConstructor{cong↑}\AgdaSpace{}%
\AgdaSymbol{:}\AgdaSpace{}%
\AgdaGeneralizable{n}\AgdaSpace{}%
\AgdaOperator{\AgdaDatatype{VRel,}}\AgdaSpace{}%
\AgdaGeneralizable{w}\AgdaSpace{}%
\AgdaOperator{\AgdaDatatype{===}}\AgdaSpace{}%
\AgdaGeneralizable{v}\AgdaSpace{}%
\AgdaSymbol{→}\AgdaSpace{}%
\AgdaSymbol{(}\AgdaInductiveConstructor{₁₊}\AgdaSpace{}%
\AgdaGeneralizable{n}\AgdaSymbol{)}\AgdaSpace{}%
\AgdaOperator{\AgdaDatatype{VRel,}}\AgdaSpace{}%
\AgdaGeneralizable{w}\AgdaSpace{}%
\AgdaOperator{\AgdaFunction{↑}}\AgdaSpace{}%
\AgdaOperator{\AgdaDatatype{===}}\AgdaSpace{}%
\AgdaGeneralizable{v}\AgdaSpace{}%
\AgdaOperator{\AgdaFunction{↑}}\<%
\\
%
\>[2]\AgdaInductiveConstructor{comm₁}\AgdaSpace{}%
\AgdaSymbol{:}\AgdaSpace{}%
\AgdaSymbol{(}\AgdaBound{h}\AgdaSpace{}%
\AgdaSymbol{:}\AgdaSpace{}%
\AgdaDatatype{Gate}\AgdaSpace{}%
\AgdaNumber{1}\AgdaSymbol{)}\AgdaSpace{}%
\AgdaSymbol{(}\AgdaBound{g}\AgdaSpace{}%
\AgdaSymbol{:}\AgdaSpace{}%
\AgdaDatatype{Gen}\AgdaSpace{}%
\AgdaGeneralizable{n}\AgdaSymbol{)}\AgdaSpace{}%
\AgdaSymbol{→}\AgdaSpace{}%
\AgdaSymbol{(}\AgdaInductiveConstructor{₁₊}\AgdaSpace{}%
\AgdaGeneralizable{n}\AgdaSymbol{)}\AgdaSpace{}%
\AgdaOperator{\AgdaDatatype{VRel,}}\AgdaSpace{}%
\AgdaOperator{\AgdaInductiveConstructor{[}}\AgdaSpace{}%
\AgdaBound{g}\AgdaSpace{}%
\AgdaOperator{\AgdaInductiveConstructor{↥}}\AgdaSpace{}%
\AgdaOperator{\AgdaInductiveConstructor{]ʷ}}\AgdaSpace{}%
\AgdaOperator{\AgdaInductiveConstructor{•}}\AgdaSpace{}%
\AgdaOperator{\AgdaInductiveConstructor{[}}\AgdaSpace{}%
\AgdaInductiveConstructor{gate₁}\AgdaSpace{}%
\AgdaBound{h}\AgdaSpace{}%
\AgdaOperator{\AgdaInductiveConstructor{]ʷ}}\AgdaSpace{}%
\AgdaOperator{\AgdaDatatype{===}}\AgdaSpace{}%
\AgdaOperator{\AgdaInductiveConstructor{[}}\AgdaSpace{}%
\AgdaInductiveConstructor{gate₁}\AgdaSpace{}%
\AgdaBound{h}\AgdaSpace{}%
\AgdaOperator{\AgdaInductiveConstructor{]ʷ}}\AgdaSpace{}%
\AgdaOperator{\AgdaInductiveConstructor{•}}\AgdaSpace{}%
\AgdaOperator{\AgdaInductiveConstructor{[}}\AgdaSpace{}%
\AgdaBound{g}\AgdaSpace{}%
\AgdaOperator{\AgdaInductiveConstructor{↥}}\AgdaSpace{}%
\AgdaOperator{\AgdaInductiveConstructor{]ʷ}}\<%
\\
%
\>[2]\AgdaInductiveConstructor{comm₂}\AgdaSpace{}%
\AgdaSymbol{:}\AgdaSpace{}%
\AgdaSymbol{(}\AgdaBound{h}\AgdaSpace{}%
\AgdaSymbol{:}\AgdaSpace{}%
\AgdaDatatype{Gate}\AgdaSpace{}%
\AgdaNumber{2}\AgdaSymbol{)}\AgdaSpace{}%
\AgdaSymbol{(}\AgdaBound{g}\AgdaSpace{}%
\AgdaSymbol{:}\AgdaSpace{}%
\AgdaDatatype{Gen}\AgdaSpace{}%
\AgdaGeneralizable{n}\AgdaSymbol{)}\AgdaSpace{}%
\AgdaSymbol{→}\AgdaSpace{}%
\AgdaSymbol{(}\AgdaInductiveConstructor{₂₊}\AgdaSpace{}%
\AgdaGeneralizable{n}\AgdaSymbol{)}\AgdaSpace{}%
\AgdaOperator{\AgdaDatatype{VRel,}}\AgdaSpace{}%
\AgdaOperator{\AgdaInductiveConstructor{[}}\AgdaSpace{}%
\AgdaBound{g}\AgdaSpace{}%
\AgdaOperator{\AgdaInductiveConstructor{↥}}\AgdaSpace{}%
\AgdaOperator{\AgdaInductiveConstructor{↥}}\AgdaSpace{}%
\AgdaOperator{\AgdaInductiveConstructor{]ʷ}}\AgdaSpace{}%
\AgdaOperator{\AgdaInductiveConstructor{•}}\AgdaSpace{}%
\AgdaOperator{\AgdaInductiveConstructor{[}}\AgdaSpace{}%
\AgdaInductiveConstructor{gate₂}\AgdaSpace{}%
\AgdaBound{h}\AgdaSpace{}%
\AgdaOperator{\AgdaInductiveConstructor{]ʷ}}\AgdaSpace{}%
\AgdaOperator{\AgdaDatatype{===}}\AgdaSpace{}%
\AgdaOperator{\AgdaInductiveConstructor{[}}\AgdaSpace{}%
\AgdaInductiveConstructor{gate₂}\AgdaSpace{}%
\AgdaBound{h}\AgdaSpace{}%
\AgdaOperator{\AgdaInductiveConstructor{]ʷ}}\AgdaSpace{}%
\AgdaOperator{\AgdaInductiveConstructor{•}}\AgdaSpace{}%
\AgdaOperator{\AgdaInductiveConstructor{[}}\AgdaSpace{}%
\AgdaBound{g}\AgdaSpace{}%
\AgdaOperator{\AgdaInductiveConstructor{↥}}\AgdaSpace{}%
\AgdaOperator{\AgdaInductiveConstructor{↥}}\AgdaSpace{}%
\AgdaOperator{\AgdaInductiveConstructor{]ʷ}}\<%
\end{code}


We use $\approx$ to denote the monoid congruence generated by
$n\;\aid{VRel,{\_}==={\_}}$, i.e., the least relation containing the
axioms and closed under reflexivity, symmetry, transitivity,
congruence of \aidop{•}, associativity, and the unit laws.

\iffalse

Because $\sigma$ is its own inverse and shifts of invertible
generators are invertible, the presented monoid is a group.

\subsection{Cosets and the Staircase Normal Form}\label{app:sn-cosets}

The engine behind the normal form is coset enumeration along the chain
$S_1 < S_2 < \dots < S_n < S_{n+1}$, where $S_n$ sits inside $S_{n+1}$ as
the permutations fixing the bottom wire --- syntactically, the circuits
of the form $c\;{\uparrow}$.  Read as permutations, two circuits lie in the
same right coset of $S_n$ exactly when they bring the same wire to the
bottom position; there are $n{+}1$ choices, so $S_n$ has index $n{+}1$,
and a natural representative of each coset is the \emph{descending
staircase} of swaps that carries that wire down.  In the library a coset
is not a set of circuits but a label naming one, an element of the finite type
$\aid{C}\;n$; here the label is the length of the staircase, as a unary
numeral bounded by the width, and the \emph{section} $[\,\_\,]^{\mathsf
c}$ maps it to the staircase circuit.  The representatives on three,
two, and one wires are drawn on the right:

\begin{center}
\begin{minipage}[c]{0.30\linewidth}
\begingroup\setlength{\mathindent}{0pt}
\begin{code}%
\>[0]\AgdaKeyword{data}\AgdaSpace{}%
\AgdaDatatype{C}\AgdaSpace{}%
\AgdaSymbol{:}\AgdaSpace{}%
\AgdaDatatype{ℕ}\AgdaSpace{}%
\AgdaSymbol{→}\AgdaSpace{}%
\AgdaPrimitive{Set}\AgdaSpace{}%
\AgdaKeyword{where}\<%
\\
\>[0][@{}l@{\AgdaIndent{0}}]%
\>[2]\AgdaInductiveConstructor{ε}%
\>[6]\AgdaSymbol{:}\AgdaSpace{}%
\AgdaDatatype{C}\AgdaSpace{}%
\AgdaGeneralizable{n}\<%
\\
%
\>[2]\AgdaOperator{\AgdaInductiveConstructor{σ•\AgdaUnderscore{}}}\AgdaSpace{}%
\AgdaSymbol{:}\AgdaSpace{}%
\AgdaDatatype{C}\AgdaSpace{}%
\AgdaGeneralizable{n}\AgdaSpace{}%
\AgdaSymbol{→}\AgdaSpace{}%
\AgdaDatatype{C}\AgdaSpace{}%
\AgdaSymbol{(}\AgdaInductiveConstructor{₁₊}\AgdaSpace{}%
\AgdaGeneralizable{n}\AgdaSymbol{)}\<%
\\
\>[0]\AgdaOperator{\AgdaFunction{[\AgdaUnderscore{}]ᶜ}}\AgdaSpace{}%
\AgdaSymbol{:}\AgdaSpace{}%
\AgdaDatatype{C}\AgdaSpace{}%
\AgdaGeneralizable{n}\AgdaSpace{}%
\AgdaSymbol{→}\AgdaSpace{}%
\AgdaFunction{Circuit}\AgdaSpace{}%
\AgdaSymbol{(}\AgdaInductiveConstructor{₁₊}\AgdaSpace{}%
\AgdaGeneralizable{n}\AgdaSymbol{)}\<%
\\
\>[0]\AgdaOperator{\AgdaFunction{[}}\AgdaSpace{}%
\AgdaInductiveConstructor{ε}\AgdaSpace{}%
\AgdaOperator{\AgdaFunction{]ᶜ}}%
\>[11]\AgdaSymbol{=}\AgdaSpace{}%
\AgdaInductiveConstructor{ε}\<%
\\
\>[0]\AgdaOperator{\AgdaFunction{[}}\AgdaSpace{}%
\AgdaOperator{\AgdaInductiveConstructor{σ•}}\AgdaSpace{}%
\AgdaBound{c}\AgdaSpace{}%
\AgdaOperator{\AgdaFunction{]ᶜ}}%
\>[11]\AgdaSymbol{=}\AgdaSpace{}%
\AgdaFunction{σ}\AgdaSpace{}%
\AgdaOperator{\AgdaInductiveConstructor{•}}\AgdaSpace{}%
\AgdaSymbol{(}\AgdaOperator{\AgdaFunction{[}}\AgdaSpace{}%
\AgdaBound{c}\AgdaSpace{}%
\AgdaOperator{\AgdaFunction{]ᶜ}}\AgdaSpace{}%
\AgdaOperator{\AgdaFunction{↑}}\AgdaSymbol{)}\<%
\end{code}
\endgroup
\end{minipage}\hspace{2.5em}%
\begin{minipage}[c]{0.42\linewidth}\centering
\begin{tabular}{@{}l@{\;\;}ccc@{}}
  $\aid{C}\;2$: & \snfig[0.35]{c2-0} & \snfig[0.35]{c2-1} & \snfig[0.35]{c2-2}\\
  & $\varepsilon$ & $\sigma{\bullet}\:\varepsilon$ &
    $\sigma{\bullet}\:\sigma{\bullet}\:\varepsilon$\\[4pt]
  $\aid{C}\;1$, $\aid{C}\;0$: & \snfig[0.35]{c1-0} & \snfig[0.35]{c1-1} & \snfig[0.35]{c0-0}\\
  & $\varepsilon$ & $\sigma{\bullet}\:\varepsilon$ & $\varepsilon$
\end{tabular}
\end{minipage}
\end{center}

\noindent So $\aid{C}\;n$ has exactly the $n{+}1$ elements drawn above
for $n = 2, 1, 0$, matching the index of $S_n$ in $S_{n+1}$.  Why this yields
a normal form is the classical argument of computational group
theory~\cite{sims1994computation,holt2005handbook}:
every $f \in S_{n+1}$ lies in one coset, say with representative $c_n$,
so $f = f' \cdot c_n$ with $f' \in S_n$, and recursing on $f'$ writes $f$
uniquely as $c_1 c_2 \cdots c_n$ with $c_k \in \aid{C}\;k$ --- a numeral
in the \emph{factorial number system}, one digit per level, since
$|S_{n+1}| = (n{+}1) \cdot n \cdots 2$.

\emph{Computing} this factorization is the job of the coset table.
Given a coset, represented by its staircase, and one generator, the
table says where their product lands --- the updated coset --- and what
is left over on the wires above --- the residual circuit --- so that
staircase followed by generator equals residual (shifted up) followed by
updated staircase.
Pushing $\sigma$ into a staircase can lengthen it, cancel with it ($\sigma^2 =
\varepsilon$), or pass it through and shift it up (Yang--Baxter), or
commute past; for a generator that does not touch the bottom wire we can
recurse.  The residual is empty or a single swap here, but in richer
gate sets it can be a longer circuit, which is why the table returns a
circuit on one wire fewer rather than a generator
(\Cref{fig:sn-pushing} draws the four cases):

\begin{center}
\begin{minipage}[c]{0.65\linewidth}
\begingroup\setlength{\mathindent}{0pt}
\begin{code}%
\>[0]\AgdaFunction{ract}\AgdaSpace{}%
\AgdaSymbol{:}\AgdaSpace{}%
\AgdaDatatype{C}\AgdaSpace{}%
\AgdaSymbol{(}\AgdaInductiveConstructor{₁₊}\AgdaSpace{}%
\AgdaGeneralizable{n}\AgdaSymbol{)}\AgdaSpace{}%
\AgdaSymbol{→}\AgdaSpace{}%
\AgdaDatatype{Gen}\AgdaSpace{}%
\AgdaSymbol{(}\AgdaInductiveConstructor{₂₊}\AgdaSpace{}%
\AgdaGeneralizable{n}\AgdaSymbol{)}\AgdaSpace{}%
\AgdaSymbol{→}\AgdaSpace{}%
\AgdaFunction{Circuit}\AgdaSpace{}%
\AgdaSymbol{(}\AgdaInductiveConstructor{₁₊}\AgdaSpace{}%
\AgdaGeneralizable{n}\AgdaSymbol{)}\AgdaSpace{}%
\AgdaOperator{\AgdaFunction{×}}\AgdaSpace{}%
\AgdaDatatype{C}\AgdaSpace{}%
\AgdaSymbol{(}\AgdaInductiveConstructor{₁₊}\AgdaSpace{}%
\AgdaGeneralizable{n}\AgdaSymbol{)}\<%
\\
\>[0]\AgdaFunction{ract}\AgdaSpace{}%
\AgdaInductiveConstructor{ε}\AgdaSpace{}%
\AgdaInductiveConstructor{σ-gen}\AgdaSpace{}%
\AgdaSymbol{=}\AgdaSpace{}%
\AgdaInductiveConstructor{ε}\AgdaSpace{}%
\AgdaOperator{\AgdaInductiveConstructor{,}}\AgdaSpace{}%
\AgdaOperator{\AgdaInductiveConstructor{σ•}}\AgdaSpace{}%
\AgdaInductiveConstructor{ε}%
\>[36]\AgdaComment{--\ grows}\<%
\\
\>[0]\AgdaFunction{ract}\AgdaSpace{}%
\AgdaSymbol{(}\AgdaOperator{\AgdaInductiveConstructor{σ•}}\AgdaSpace{}%
\AgdaInductiveConstructor{ε}\AgdaSymbol{)}\AgdaSpace{}%
\AgdaInductiveConstructor{σ-gen}\AgdaSpace{}%
\AgdaSymbol{=}\AgdaSpace{}%
\AgdaInductiveConstructor{ε}\AgdaSpace{}%
\AgdaOperator{\AgdaInductiveConstructor{,}}\AgdaSpace{}%
\AgdaInductiveConstructor{ε}%
\>[36]\AgdaComment{--\ cancels\ (σ²\ =\ ε)}\<%
\\
\>[0]\AgdaFunction{ract}\AgdaSpace{}%
\AgdaSymbol{(}\AgdaOperator{\AgdaInductiveConstructor{σ•}}\AgdaSpace{}%
\AgdaOperator{\AgdaInductiveConstructor{σ•}}\AgdaSpace{}%
\AgdaBound{c}\AgdaSymbol{)}\AgdaSpace{}%
\AgdaInductiveConstructor{σ-gen}\AgdaSpace{}%
\AgdaSymbol{=}\AgdaSpace{}%
\AgdaFunction{σ}\AgdaSpace{}%
\AgdaOperator{\AgdaInductiveConstructor{,}}\AgdaSpace{}%
\AgdaOperator{\AgdaInductiveConstructor{σ•}}\AgdaSpace{}%
\AgdaOperator{\AgdaInductiveConstructor{σ•}}\AgdaSpace{}%
\AgdaBound{c}%
\>[36]\AgdaComment{--\ passes\ (Yang-Baxter)}\<%
\\
\>[0]\AgdaFunction{ract}\AgdaSpace{}%
\AgdaInductiveConstructor{ε}\AgdaSpace{}%
\AgdaSymbol{(}\AgdaBound{g}\AgdaSpace{}%
\AgdaOperator{\AgdaInductiveConstructor{↥}}\AgdaSymbol{)}\AgdaSpace{}%
\AgdaSymbol{=}\AgdaSpace{}%
\AgdaOperator{\AgdaInductiveConstructor{[}}\AgdaSpace{}%
\AgdaBound{g}\AgdaSpace{}%
\AgdaOperator{\AgdaInductiveConstructor{]ʷ}}\AgdaSpace{}%
\AgdaOperator{\AgdaInductiveConstructor{,}}\AgdaSpace{}%
\AgdaInductiveConstructor{ε}%
\>[36]\AgdaComment{--\ commutes\ past}\<%
\\
\>[0]\AgdaFunction{ract}\AgdaSpace{}%
\AgdaSymbol{(}\AgdaOperator{\AgdaInductiveConstructor{σ•}}\AgdaSpace{}%
\AgdaBound{c}\AgdaSymbol{)}\AgdaSpace{}%
\AgdaSymbol{(}\AgdaBound{g}\AgdaSpace{}%
\AgdaOperator{\AgdaInductiveConstructor{↥}}\AgdaSymbol{)}\AgdaSpace{}%
\AgdaSymbol{=}\AgdaSpace{}%
\AgdaField{proj₁}\AgdaSpace{}%
\AgdaSymbol{(}\AgdaFunction{ract}\AgdaSpace{}%
\AgdaBound{c}\AgdaSpace{}%
\AgdaBound{g}\AgdaSymbol{)}\AgdaSpace{}%
\AgdaOperator{\AgdaFunction{↑}}\AgdaSpace{}%
\AgdaOperator{\AgdaInductiveConstructor{,}}\AgdaSpace{}%
\AgdaOperator{\AgdaInductiveConstructor{σ•}}\AgdaSpace{}%
\AgdaSymbol{(}\AgdaField{proj₂}\AgdaSpace{}%
\AgdaSymbol{(}\AgdaFunction{ract}\AgdaSpace{}%
\AgdaBound{c}\AgdaSpace{}%
\AgdaBound{g}\AgdaSymbol{))}%
\>[64]\AgdaComment{--\ recurses}\<%
\end{code}
\endgroup
\end{minipage}\hspace{0.5em}%
\begin{minipage}[c]{0.30\linewidth}\raggedright
  \snfigT[0.5]{ract-l}\;\scalebox{0.7}{$\approx$}\;\snfigT[0.5]{ract-r}
\end{minipage}
\end{center}

\begin{figure}[!htbp]
\centering
\begin{tabular}{@{}r@{\,}c@{\,}l@{\;\;}l@{\qquad}r@{\,}c@{\,}l@{\;\;}l@{}}
  \snfigT[0.4]{push4-l} & \scalebox{0.7}{$\rightarrow$} & \snfigT[0.4]{push4-r} &
    \small grows &
  \snfigT[0.4]{push3-l} & \scalebox{0.7}{$\rightarrow$} & \snfigT[0.4]{push3-r} &
    \small cancels ($\sigma^2 = \varepsilon$)\\[4pt]
  \snfigT[0.4]{push2-l} & \scalebox{0.7}{$\rightarrow$} & \snfigT[0.4]{push2-r} &
    \small passes (Yang--Baxter) &
  \snfigT[0.4]{push1-l} & \scalebox{0.7}{$\rightarrow$} & \snfigT[0.4]{push1-r} &
    \small commutes past
\end{tabular}
\caption{Pushing a swap through a staircase: the four cases of the coset
table \aid{ract} as circuit rewrites.}
\label{fig:sn-pushing}
\Description{Four cases for pushing a swap through a staircase:
the staircase grows, two swaps cancel, a swap passes through by the
Yang--Baxter relation, or a disjoint swap commutes past. The labels mark
the old and new coset representatives and residual circuits.}
\end{figure}

\medskip\noindent
\aid{σ-gen} abbreviates $\aid{gate₂}\;\aid{σ-gate}$.  The last clause is the recursion: a
shifted generator meets a staircase by meeting its tail one wire down.
Each clause is certified by the \emph{soundness law}, pictured on the
right and proved once per table by a short equational derivation from
the axioms:

\begin{code}%
\>[0]\AgdaFunction{ract-sound}\AgdaSpace{}%
\AgdaSymbol{:}\AgdaSpace{}%
\AgdaSymbol{∀}\AgdaSpace{}%
\AgdaBound{c}\AgdaSpace{}%
\AgdaBound{b}\AgdaSpace{}%
\AgdaSymbol{→}\AgdaSpace{}%
\AgdaKeyword{let}\AgdaSpace{}%
\AgdaSymbol{(}\AgdaBound{b'}\AgdaSpace{}%
\AgdaOperator{\AgdaInductiveConstructor{,}}\AgdaSpace{}%
\AgdaBound{c'}\AgdaSymbol{)}\AgdaSpace{}%
\AgdaSymbol{=}\AgdaSpace{}%
\AgdaFunction{ract}\AgdaSpace{}%
\AgdaBound{c}\AgdaSpace{}%
\AgdaBound{b}\AgdaSpace{}%
\AgdaKeyword{in}\AgdaSpace{}%
\AgdaOperator{\AgdaFunction{[}}\AgdaSpace{}%
\AgdaBound{c}\AgdaSpace{}%
\AgdaOperator{\AgdaFunction{]ᶜ}}\AgdaSpace{}%
\AgdaOperator{\AgdaInductiveConstructor{•}}\AgdaSpace{}%
\AgdaOperator{\AgdaInductiveConstructor{[}}\AgdaSpace{}%
\AgdaBound{b}\AgdaSpace{}%
\AgdaOperator{\AgdaInductiveConstructor{]ʷ}}\AgdaSpace{}%
\AgdaOperator{\AgdaPostulate{≈}}\AgdaSpace{}%
\AgdaBound{b'}\AgdaSpace{}%
\AgdaOperator{\AgdaFunction{↑}}\AgdaSpace{}%
\AgdaOperator{\AgdaInductiveConstructor{•}}\AgdaSpace{}%
\AgdaOperator{\AgdaFunction{[}}\AgdaSpace{}%
\AgdaBound{c'}\AgdaSpace{}%
\AgdaOperator{\AgdaFunction{]ᶜ}}\<%
\end{code}
\begin{code}[hide]%
\>[0]\AgdaFunction{ract-sound}\AgdaSpace{}%
\AgdaSymbol{=}\AgdaSpace{}%
\AgdaPostulate{proof-elided}\<%
\end{code}

The stateful word traversal \aid{{\_}ᵗ} extends
\aid{ract} to whole circuits, threading the coset from left to right and
concatenating the residuals; run from the identity coset, it maps an
$(n{+}2)$-wire circuit to an $(n{+}1)$-wire circuit and a coset --- one
level of normalization.

\emph{Iterating} the level is the coset tower, whose carrier grows by
one digit per level:

\begin{center}
\begin{minipage}[c]{0.56\linewidth}
\begingroup\setlength{\mathindent}{0pt}
\begin{code}%
\>[0]\AgdaFunction{NF}\AgdaSpace{}%
\AgdaSymbol{:}\AgdaSpace{}%
\AgdaDatatype{ℕ}\AgdaSpace{}%
\AgdaSymbol{→}\AgdaSpace{}%
\AgdaPrimitive{Set}\<%
\\
\>[0]\AgdaFunction{NF}\AgdaSpace{}%
\AgdaNumber{0}%
\>[11]\AgdaSymbol{=}\AgdaSpace{}%
\AgdaRecord{⊤}\<%
\\
\>[0]\AgdaFunction{NF}\AgdaSpace{}%
\AgdaSymbol{(}\AgdaInductiveConstructor{suc}\AgdaSpace{}%
\AgdaBound{k}\AgdaSymbol{)}\AgdaSpace{}%
\AgdaSymbol{=}\AgdaSpace{}%
\AgdaFunction{NF}\AgdaSpace{}%
\AgdaBound{k}\AgdaSpace{}%
\AgdaOperator{\AgdaFunction{×}}\AgdaSpace{}%
\AgdaDatatype{C}\AgdaSpace{}%
\AgdaBound{k}%
\>[26]\AgdaComment{--\ ⊤\ ×\ C\ 0\ ×\ C\ 1\ ×\ ⋯\ ×\ C\ k}\<%
\end{code}
\endgroup
\end{minipage}\hspace{2.5em}%
\begin{minipage}[c]{0.16\linewidth}\centering
  \snfig[0.35]{fig3_staircase}
\end{minipage}
\end{center}

\noindent so $\aid{NF}\;n$ has exactly $n!$ elements.  The normal-form
function $\aid{nf} : \aid{Circuit}\;n \to \aid{NF}\;n$ runs the tables;
its \emph{section} $\aid{inv-nf} : \aid{NF}\;n \to \aid{Circuit}\;n$
rebuilds a circuit from its digits by concatenating the staircases,
suitably shifted, as drawn on the right for $(\mathsf{tt}, c_0, c_1,
c_2, c_3) \in \aid{NF}\;4$, the digit $c_0 \in \aid{C}\;0$ being the
empty staircase --- one staircase per level, each shifted up to the
wires of its level: a ``staircase of staircases''.

\fi

